\documentclass{article}
\title{Superscripts and Subscripts}
\author{ }
\date{}

\begin{document}
\maketitle

\section*{1. What are superscripts and subscripts?}
A superscript is a small character placed above the normal line, often used for powers. A subscript is a small character placed below the line, often used for labels or sequence numbers.

\section*{2. Superscripts}
\textbf{LaTeX source code:}
\begin{verbatim}
\[
x^2
\]
\end{verbatim}
\textbf{Expected compiled output:}
\[
x^2
\]
\textbf{Explanation:} The caret symbol \verb|^| creates a superscript. Here, 2 is the power of \(x\).

\section*{3. Subscripts}
\textbf{LaTeX source code:}
\begin{verbatim}
\[
a_1
\]
\end{verbatim}
\textbf{Expected compiled output:}
\[
a_1
\]
\textbf{Explanation:} The underscore \verb|_| creates a subscript. Here, 1 is written below \(a\).

\section*{4. Grouping multiple characters}
Use braces when the superscript or subscript contains more than one character.

\textbf{LaTeX source code:}
\begin{verbatim}
\[
x^{10} \qquad a_{12}
\]
\end{verbatim}
\textbf{Expected compiled output:}
\[
x^{10} \qquad a_{12}
\]

\section*{5. Practice}
Write \(y^3\), \(b_2\), and \(m_{15}\) in LaTeX.
\end{document}
